\documentclass[11pt,a4paper]{article}
\usepackage[T1]{fontenc}
\usepackage[utf8]{inputenc}
\usepackage[ngerman]{babel}
\usepackage{lmodern}
\usepackage{microtype}
\usepackage{amssymb}
\usepackage[margin=2.2cm]{geometry}
\usepackage{array}
\usepackage{booktabs}
\usepackage{longtable}
\usepackage{enumitem}
\usepackage{xcolor}
\usepackage[hidelinks]{hyperref}
\setlength{\parindent}{0pt}
\setlength{\parskip}{0.45em}
\setlist[itemize]{leftmargin=*,itemsep=0.15em,topsep=0.2em}
\setlist[enumerate]{leftmargin=*,itemsep=0.25em,topsep=0.2em}
\setlength{\emergencystretch}{3em}
\newcommand{\todoitem}{\(\square\)\enspace}
\newcommand{\datenstand}{25. November 2026}
\newcommand{\abgabe}{16. Februar 2027}
\newcommand{\gespraech}{27. November 2026}

\title{Zeit- und Arbeitsplan zur Bachelorarbeit\\
\large Literaturrecherche, Experiment und Schreiben}
\author{Marco Garcia Jung}
\date{Entwurf vom \datenstand}

\begin{document}
\maketitle

\textbf{Status: Entwurf zur persönlichen Freigabe und Anpassung.}
Dieser Plan ist keine Änderung des abgestimmten Exposés. Die Daten sind
interne Zieltermine, keine zusätzlichen Abgabefristen. Tatsächliche
Arbeitszeiten, Feiertage und Rückmeldungen der Betreuung können Anpassungen
erfordern.

\begin{tabular}{@{}p{4.0cm}p{11.3cm}@{}}
\toprule
Abgabe & \abgabe{}; ab 25. November noch 83 Kalendertage \\
Betreuungsgespräch & \gespraech{} \\
Umfang & ungefähr 40 Seiten als vorläufige Annahme; Vorgabe noch klären \\
Versuch & TabPFN-3.5 und LightGBM; MCAR, MAR und MNAR vorgesehen \\
Arbeitsorte & Hauptversuch am PC; Recherche und Schreiben auch am Laptop \\
\bottomrule
\end{tabular}

\section{Wo stehe ich?}

\begin{itemize}
  \item \textbf{Erledigt:} Exposé und Forschungsfragen stehen, die Arbeit ist
  angemeldet. Literaturbestand, Suchprotokoll, Konzeptmatrix und
  LaTeX-Kapitelgerüst sind vorhanden.
  \item \textbf{Literatur:} Eine neue, nachvollziehbar dokumentierte Suche
  wurde noch nicht begonnen. Die vorbefüllten Matrixzeilen sind am Volltext
  zu prüfen. Der vorhandene Bestand bleibt erhalten, wird aber nicht
  nachträglich als Treffer einer nicht durchgeführten Suche ausgegeben.
  \item \textbf{Technischer Stand:} Ein WDBC-Pilot mit TabPFN-3.5 und
  LightGBM und ein MCAR-Laufzeittest auf Spambase liegen vor. Der PC besitzt
  eine NVIDIA RTX 5070 Ti mit 16 GB. Die Tests belegen technische Teilaspekte,
  nicht die Ergebnisse oder vollständige Machbarkeit des Hauptversuchs.
  \item \textbf{Noch offen:} Regeln und Umfang von MAR/MNAR, endgültige
  Datensätze, Splits, Wiederholungen, Tuningbudget, Auswertung und die
  Implementierung des Hauptversuchs. Der Versuch wird erst nach Festlegung
  und Prüfung des Designs gestartet.
\end{itemize}

\textbf{Zusammenhang:} Die Literaturanalyse begründet das Versuchsdesign;
Implementierung und Tests können beginnen, während relevante Literatur
weiter geprüft wird. Schreiben läuft parallel, Einleitung und Fazit werden
am Ende fertiggestellt.

\section{Meilensteine bis zur Abgabe}

\small
\begin{longtable}{@{}p{2.6cm}p{6.5cm}p{6.0cm}@{}}
\toprule
\textbf{Zeitraum} & \textbf{Arbeitspaket} & \textbf{Ergebnis am Ende} \\
\midrule
\endfirsthead
\toprule
\textbf{Zeitraum} & \textbf{Arbeitspaket} & \textbf{Ergebnis am Ende} \\
\midrule
\endhead
25.--27.11. & Plan freigeben, Suchprotokoll vorbereiten, erste dokumentierte
Suchläufe und Sichtung; Betreuungsgespräch. & Tatsächlicher Zwischenstand und
schriftlich festgehaltene Rückmeldungen. \\
\addlinespace
28.11.--04.12. & Suchblöcke S1/S2 bearbeiten; Suchläufe protokollieren,
Treffer in Zotero erfassen, Dubletten und Titel/Abstracts sichten. &
Dokumentierte Suchläufe und priorisierte Volltextliste. \\
\addlinespace
05.--11.12. & S3/S4, Vorwärts- und Rückwärtssuche; Volltexte gezielt prüfen,
Konzeptmatrix fortschreiben. & Vergleichbare, mit Fundstellen belegte
Matrixeinträge; offene Fragen sichtbar. \\
\addlinespace
12.--18.12. & Literatur konzeptzentriert zusammenführen, Grenzen und
Untersuchungsbedarf festhalten; Kapitel 3 als Rohfassung. &
Rohfassung FF1 und begründete Vorschläge für das Experiment. \\
\addlinespace
19.--27.12. & Datenlader, Gruppensplits, Maskierung und Behandlungen
implementieren; erste Tests. Parallel Kapitel 2/4 beginnen. &
Einzeln testbare Komponenten und dokumentierte offene Designpunkte. \\
\addlinespace
28.12.--03.01. & Mit Betreuung offene Designfragen abschließen; Splits,
Masken, Tuning, Auswertung, Seeds und Umgebung festlegen. Vollständigen
Probedurchlauf und Fehlerbehandlung testen. &
Vorab festgelegtes Versuchsprotokoll und geprüfte End-to-End-Pipeline. \\
\addlinespace
04.--10.01. & Hauptversuch am PC durchführen; Konfiguration,
Versionen, Checkpoint, Seeds und Einzelresultate speichern. &
Gesicherte Ergebnisdateien für die tatsächlich durchgeführten Bedingungen. \\
\addlinespace
11.--17.01. & Ergebnisse pro Datensatz auswerten, Abbildungen und Tabellen
erstellen; Methodik und Ergebniskapitel schreiben. &
Kapitel 4 und 5 als prüfbare Rohfassung. \\
\addlinespace
18.--24.01. & Grundlagen vervollständigen; Ergebnisse, Literatur und
Limitationen in der Diskussion zusammenführen. &
Rohfassungen von Kapitel 2 und 6. \\
\addlinespace
25.01.--01.02. & Einleitung, Fazit und Abstract schreiben;
Forschungsfragen und roten Faden über alle Kapitel abgleichen. &
Vollständige Rohfassung bis spätestens 01.02. \\
\addlinespace
02.--08.02. & Inhalt und Argumentation überarbeiten; Quellenstellen,
Zahlen, Abbildungen und Literaturverzeichnis prüfen. &
Inhaltlich überprüfte Fassung; offene Punkte dokumentiert. \\
\addlinespace
09.--15.02. & Sprache, Format, Anhang und KI-Angaben prüfen; PDF-Build,
Abgabevorgaben und gegebenenfalls Druck oder Upload vorbereiten. &
Abgabefertige Fassung und Puffer für technische Probleme. \\
\addlinespace
16.02. & Abgabe gemäß den geltenden Vorgaben. & Arbeit eingereicht. \\
\bottomrule
\end{longtable}
\normalsize

\textbf{Puffer:} Die zwei Wochen ab 2. Februar sind für Überarbeitung und
Abgabe vorgesehen. Feiertage und persönliche Verpflichtungen insbesondere
Ende Dezember früh eintragen; freie Tage nicht als sichere Arbeitszeit
verplanen. Bei Verzögerung zusätzliche Recherche gezielt eingrenzen und die
Betreuung informieren, statt ungetestete Versuche durchzuführen oder Quellen-
und Fehlerprüfung zu streichen.

\section{Nach Freigabe: intensiver Start am 25. November}

Die folgenden Blöcke sind ein \textbf{Vorschlag}, keine Annahme über deine
verfügbare Zeit. Zusammen etwa 4,5--6 Stunden zuzüglich Pausen. Wenn weniger
Zeit verfügbar ist, Blöcke 1--3 zuerst bearbeiten und den Rest verschieben.

\begin{enumerate}
  \item \textbf{Vorbereiten, 20--30 Minuten.} In Zotero Sammlungen für
  vorhandenen Bestand, neue Treffer und zu prüfende Volltexte anlegen.
  Suchprotokoll öffnen. Vor Beginn Abschnitt 2 an den beschlossenen
  Neustart anpassen: Fundwege tatsächlich dokumentieren; Bestand nicht
  rückwirkend als Suchtreffer zählen.
  \item \textbf{Suchplan prüfen, 30--45 Minuten.} Forschungsfrage und
  S1--S4 durchgehen. Die bestehenden Ein- und Ausschlusskriterien sind
  Vorschläge; Grundlagen und reine Rekonstruktionsstudien gesondert als
  Kontext führen, sofern sie für das Verständnis gebraucht werden. Suchsyntax,
  Felder, Filter und Sichtungsgrenzen vor jedem Suchlauf festhalten.
  \item \textbf{Erste Suchläufe, 90--120 Minuten.} In Google Scholar zunächst
  \texttt{"TabPFN" "missing values"} suchen und als Teil-Suche von S1
  protokollieren. Anschließend eine weitere passende S1-Variante aus dem
  Suchprotokoll ausführen. Datum, genaue Eingabe, Sortierung, Filter,
  geschätzte Trefferzahl und tatsächlich gesichteten Bereich eintragen.
  Metadaten des festgelegten Sichtungsbereichs soweit möglich sichern.
  \item \textbf{Sichten, 60--90 Minuten.} Dubletten und unterschiedliche
  Fassungen prüfen; Titel und Abstracts lesen. Relevante und unklare Treffer
  zur Volltextprüfung vormerken. Anzahl und Fundwege nachvollziehbar
  dokumentieren; keine ausführliche Notiz für jeden Treffer.
  \item \textbf{Erste Volltextprüfung, 45--60 Minuten.} Eine relevante
  Arbeit auswählen, Methoden und Ergebnisse gezielt lesen und ihre
  Konzeptmatrixzeile mit Abschnitts- oder Seitenangaben ergänzen.
  Ungeprüfte Angaben mit \texttt{?} oder entsprechendem Status stehen lassen.
\end{enumerate}

\textbf{Tagesabschluss:} Such-Log und Zotero-Status prüfen; offene Punkte und
nächsten Suchlauf notieren. Es gibt keine Zielzahl von 100 Papern und keine
Pflicht, jede Einleitung zu lesen.

\section{26. und 27. November}

\subsection*{Donnerstag, 26. November}
\begin{itemize}
  \item \todoitem S1/S2 fortsetzen; nach Möglichkeit Scopus mit Uni-Zugang
  nutzen und die exakte datenbankspezifische Eingabe dokumentieren.
  \item \todoitem Eine Rückwärts- und eine Vorwärtssuche zu einer
  einschlägigen Arbeit durchführen und im Suchprotokoll festhalten.
  \item \todoitem Eine weitere relevante Arbeit am Volltext prüfen oder
  vorhandene Matrixangaben gegen ihre Fundstellen abgleichen.
  \item \todoitem Zwischenstand und Fragen für den Termin auf einer
  überschaubaren Liste sammeln.
\end{itemize}

\subsection*{Freitag, 27. November: Betreuungsgespräch}
\begin{itemize}
  \item \todoitem Zeitplan, tatsächliche Suchläufe und eine beispielhaft
  geprüfte Matrixzeile mitnehmen; unvollständige Schritte offen benennen.
  \item \todoitem Umfang der drei Fehlmechanismen, konkrete Datenbasis,
  Wiederholungen und Anforderungen an statistische Auswertung besprechen.
  \item \todoitem Verbindlichen Seitenumfang, IEEE-Stil trotz LNCS-Vorlage,
  Form der Suchdokumentation und KI-Angaben klären.
  \item \todoitem Ergebnisse nach dem Gespräch im Projektstand eintragen und
  Zeitplan bei Bedarf gemeinsam anpassen.
\end{itemize}

\section{Recherche: geplanter Datenbankrahmen}

Als Arbeitsgrundlage dienen die bereits im Suchprotokoll genannten
\textbf{Google Scholar, arXiv, Scopus, IEEE Xplore und ACM Digital Library}
sowie Vorwärts- und Rückwärtssuche. Zugang und Suchsyntax je Plattform
prüfen; Suchläufe und tatsächlich gesichtete Bereiche einzeln dokumentieren.
Das Exposé nennt ausdrücklich Scholar, arXiv und Zitationssuche; die
weiteren Datenbanken stehen im Suchprotokoll. Die Dokumentation wird nach
Freigabe und gegebenenfalls Rücksprache inhaltlich konsistent gehalten;
das Exposé wird nicht ohne Rücksprache geändert.

Der geplante Charakter bleibt eine \textbf{fokussierte dokumentierte
Literaturanalyse}, kein automatisch vollständiges systematisches Review.
Auswahlzahlen und Ausschlussgründe nachvollziehbar festhalten; die konkrete
Darstellung nach PRISMA mit der Betreuung klären. Aus geschätzten Scholar-
Trefferzahlen und begrenzter Sichtung keine exakten Gesamtzahlen ableiten.

\section{Vor dem Hauptversuch: verbindliche Prüfpunkte}

\begin{itemize}
  \item \todoitem Auswahl der Datensätze und identische Eingaben als Gruppen
  vor Maskierung und datenabhängiger Vorverarbeitung prüfen.
  \item \todoitem MCAR-, MAR- und MNAR-Regeln einschließlich Abhängigkeiten,
  Ziel- und realisierter Fehlraten, betroffener Partitionen und Seeds
  festlegen; Labels nicht maskieren.
  \item \todoitem Splits vor Maskierung; gleiche Splits und Masken für beide
  Modelle und alle Behandlungen; Median nur aus beobachteten Trainingswerten.
  \item \todoitem Leere Trainingsspalten, Indikatoren, LightGBM-Tuning,
  Validierung und Sensitivitätsanalyse vorab regeln.
  \item \todoitem Balanced Accuracy, Log-Loss und gepaarte Differenz
  \(\Delta-\Delta_0\) sowie Aggregation festlegen; Wiederholungen nicht als
  unabhängige Datensätze behandeln.
  \item \todoitem Paketversionen, Modellkonfiguration, Checkpoint und Seeds
  festhalten; Einzelergebnisse und realisierte Fehlanteile speichern.
  \item \todoitem Gesamten Ablauf mit kleinen Daten testen. Der Spambase-Test
  nutzte den Speichermodus \texttt{memory\_saving\_mode=True}; dies ist für die Hauptpipeline
  gesondert zu prüfen und zu dokumentieren.
\end{itemize}

\section{So passe ich den Zeitplan an}

Die bearbeitbare Quelle heißt \texttt{writing/Zeitplan.tex}. Daten und Texte
können dort direkt verändert werden; für das PDF im Projektverzeichnis:

\begin{verbatim}
Set-Location writing
pdflatex -interaction=nonstopmode -halt-on-error Zeitplan.tex
pdflatex -interaction=nonstopmode -halt-on-error Zeitplan.tex
\end{verbatim}

Nach Änderungen an tatsächlichen Entscheidungen oder Fristen zusätzlich
\texttt{PROJEKTSTAND.md} aktualisieren. Ein nicht erreichter interner
Meilenstein wird neu geplant; vergangene Schritte nicht als erledigt
kennzeichnen, wenn sie nicht stattgefunden haben.

\end{document}
